\section{AHSME 1957}
        \begin{enumerate}
        	\item (\textit{AHSME 1957}) The number of distinct lines representing the altitudes, medians, and interior angle bisectors of a triangle that is isosceles, but not equilateral, is: 


 $\textbf{(A)}\ 9\qquad \textbf{(B)}\ 7\qquad \textbf{(C)}\ 6\qquad \textbf{(D)}\ 5\qquad \textbf{(E)}\ 3$


 

	\item (\textit{AHSME 1957}) In the equation $2x^2 - hx + 2k = 0$, the sum of the roots is $4$ and the product of the roots is $-3$. 
Then $h$ and $k$ have the values, respectively: 


 $\textbf{(A)}\ 8\text{ and }{-6} \qquad  \textbf{(B)}\ 4\text{ and }{-3}\qquad  \textbf{(C)}\ {-3}\text{ and }4\qquad \textbf{(D)}\ {-3}\text{ and }8\qquad \textbf{(E)}\ 8\text{ and }{-3}$


 

	\item (\textit{AHSME 1957}) The simplest form of $1 - \frac{1}{1 + \frac{a}{1 - a}}$ is: 


 $\textbf{(A)}\ {a}\text{ if }{a\not= 0} \qquad \textbf{(B)}\ 1\qquad  \textbf{(C)}\ {a}\text{ if }{a\not=-1}\qquad \textbf{(D)}\ {1-a}\text{ with not restriction on }{a}\qquad \textbf{(E)}\ {a}\text{ if }{a\not= 1}$


 

	\item (\textit{AHSME 1957}) The first step in finding the product $(3x + 2)(x - 5)$ by use of the distributive property in the form $a(b + c) = ab + ac$ is: 


 $\textbf{(A)}\ 3x^2 - 13x - 10 \qquad  \textbf{(B)}\ 3x(x - 5) + 2(x - 5)\qquad \\ \textbf{(C)}\ (3x+2)x+(3x+2)(-5)\qquad \textbf{(D)}\ 3x^2-17x-10\qquad \textbf{(E)}\ 3x^2+2x-15x-10$


 

	\item (\textit{AHSME 1957}) Through the use of theorems on logarithms 

 \[\log{\frac{a}{b}} + \log{\frac{b}{c}} + \log{\frac{c}{d}} - \log{\frac{ay}{dx}}\]
can be reduced to: 


 $\textbf{(A)}\ \log{\frac{y}{x}}\qquad  \textbf{(B)}\ \log{\frac{x}{y}}\qquad  \textbf{(C)}\ 1\qquad \\ \textbf{(D)}\ 140x-24x^2+x^3\qquad \textbf{(E)}\ \text{none of these}$


 

	\item (\textit{AHSME 1957}) An open box is constructed by starting with a rectangular sheet of metal $10$ in. by $14$ in. 
and cutting a square of side $x$ inches from each corner. The resulting projections are folded up and the seams welded. 
The volume of the resulting box is: 


 $\textbf{(A)}\ 140x - 48x^2 + 4x^3 \qquad  \textbf{(B)}\ 140x + 48x^2 + 4x^3\qquad \\ \textbf{(C)}\ 140x+24x^2+x^3\qquad \textbf{(D)}\ 140x-24x^2+x^3\qquad \textbf{(E)}\ \text{none of these}$


 

	\item (\textit{AHSME 1957}) The area of a circle inscribed in an equilateral triangle is $48\pi$. The perimeter of this triangle is: 


 $\textbf{(A)}\ 72\sqrt{3} \qquad \textbf{(B)}\ 48\sqrt{3}\qquad \textbf{(C)}\ 36\qquad \textbf{(D)}\ 24\qquad \textbf{(E)}\ 72$


 

	\item (\textit{AHSME 1957}) The numbers $x,\,y,\,z$ are proportional to $2,\,3,\,5$. The sum of $x, y$, and $z$ is $100$. The number y is given by the equation $y = ax - 10$. Then a is: 


 $\textbf{(A)}\ 2 \qquad \textbf{(B)}\ \frac{3}{2}\qquad \textbf{(C)}\ 3\qquad \textbf{(D)}\ \frac{5}{2}\qquad \textbf{(E)}\ 4$


 

	\item (\textit{AHSME 1957}) The value of $x - y^{x - y}$ when $x = 2$ and $y = -2$ is: 


 $\textbf{(A)}\ -18 \qquad \textbf{(B)}\ -14\qquad \textbf{(C)}\ 14\qquad \textbf{(D)}\ 18\qquad \textbf{(E)}\ 256$


 

	\item (\textit{AHSME 1957}) The graph of $y = 2x^2 + 4x + 3$ has its: 


 $\textbf{(A)}\ \text{lowest point at } {(-1,9)}\qquad   \textbf{(B)}\ \text{lowest point at } {(1,1)}\qquad \\  \textbf{(C)}\ \text{lowest point at }{(-1,1)}\qquad \textbf{(D)}\ \text{highest point at }{(-1,9)}\qquad\\  \textbf{(E)}\ \text{highest point at }{(-1,1)}$


 

	\item (\textit{AHSME 1957}) The angle formed by the hands of a clock at $2:15$ is: 


 $\textbf{(A)}\ 30^\circ \qquad \textbf{(B)}\ 27\frac{1}{2}^\circ\qquad  \textbf{(C)}\ 157\frac{1}{2}^\circ\qquad  \textbf{(D)}\ 172\frac{1}{2}^\circ\qquad \textbf{(E)}\ \text{none of these}$


 

	\item (\textit{AHSME 1957}) Comparing the numbers $10^{-49}$ and $2\cdot 10^{-50}$ we may say: 


 $\textbf{(A)}\ \text{the first exceeds the second by }{8\cdot 10^{-1}}\qquad\\  \textbf{(B)}\ \text{the first exceeds the second by }{2\cdot 10^{-1}}\qquad\\  \textbf{(C)}\ \text{the first exceeds the second by }{8\cdot 10^{-50}}\qquad\\  \textbf{(D)}\ \text{the second is five times the first}\qquad\\  \textbf{(E)}\ \text{the first exceeds the second by }{5}$


 

	\item (\textit{AHSME 1957}) A rational number between $\sqrt{2}$ and $\sqrt{3}$ is: 


 $\textbf{(A)}\ \frac{\sqrt{2} + \sqrt{3}}{2} \qquad  \textbf{(B)}\ \frac{\sqrt{2} \cdot \sqrt{3}}{2}\qquad  \textbf{(C)}\ 1.5\qquad \textbf{(D)}\ 1.8\qquad \textbf{(E)}\ 1.4$


 

	\item (\textit{AHSME 1957}) If $y = \sqrt{x^2 - 2x + 1} + \sqrt{x^2 + 2x + 1}$, then $y$ is: 


 $\textbf{(A)}\ 2x\qquad \textbf{(B)}\ 2(x+1)\qquad\textbf{(C)}\ 0\qquad\textbf{(D)}\ |x-1|+|x+1|\qquad\textbf{(E)}\ \text{none of these}$


 

	\item (\textit{AHSME 1957}) The table below shows the distance $s$ in feet a ball rolls down an inclined plane in $t$ seconds. 


 
 \[\begin{tabular}{|c|c|c|c|c|c|c|}\hline t & 0 & 1 & 2 & 3 & 4 & 5\\ \hline s & 0 & 10 & 40 & 90 & 160 & 250\\ \hline\end{tabular}\]


 The distance $s$ for $t = 2.5$ is: 


 $\textbf{(A)}\ 45\qquad \textbf{(B)}\ 62.5\qquad \textbf{(C)}\ 70\qquad \textbf{(D)}\ 75\qquad \textbf{(E)}\ 82.5$


 

	\item (\textit{AHSME 1957}) Goldfish are sold at $15$ cents each. The rectangular coordinate graph showing the cost of $1$ to $12$ goldfish is: 


 $\textbf{(A)}\ \text{a straight line segment} \qquad \\  \textbf{(B)}\ \text{a set of horizontal parallel line segments}\qquad\\  \textbf{(C)}\ \text{a set of vertical parallel line segments}\qquad\\  \textbf{(D)}\ \text{a finite set of distinct points} \qquad\textbf{(E)}\ \text{a straight line}$


 

	\item (\textit{AHSME 1957}) A cube is made by soldering twelve $3$-inch lengths of wire properly at the vertices of the cube. 
If a fly alights at one of the vertices and then walks along the edges, the greatest distance it 
could travel before coming to any vertex a second time, without retracing any distance, is: 


 $\textbf{(A)}\ 24\text{ in.}\qquad \textbf{(B)}\ 12\text{ in.}\qquad \textbf{(C)}\ 30\text{ in.}\qquad  \textbf{(D)}\ 18\text{ in.}\qquad\textbf{(E)}\ 36\text{ in.}$


 

	\item (\textit{AHSME 1957}) Circle $O$ has diameters $AB$ and $CD$ perpendicular to each other. $AM$ is any chord intersecting $CD$ at $P$. 
Then $AP\cdot AM$ is equal to: 


 
                    \begin{figure}[H]
                        \centering
                        \includegraphics[width=0.5\textwidth]{Problem/AHSME_I_1957_18_Problem_Figure_1.png}
                    \end{figure}
                    


 $\textbf{(A)}\ AO\cdot OB \qquad \textbf{(B)}\ AO\cdot AB\qquad \\ \textbf{(C)}\ CP\cdot CD \qquad  \textbf{(D)}\ CP\cdot PD\qquad  \textbf{(E)}\ CO\cdot OP$


 

	\item (\textit{AHSME 1957}) The base of the decimal number system is ten, meaning, for example, that $123 = 1\cdot 10^2 + 2\cdot 10 + 3$. 
In the binary system, which has base two, the first five positive integers are $1,\,10,\,11,\,100,\,101$. 
The numeral $10011$ in the binary system would then be written in the decimal system as: 


 $\textbf{(A)}\ 19 \qquad \textbf{(B)}\ 40\qquad \textbf{(C)}\ 10011\qquad \textbf{(D)}\ 11\qquad \textbf{(E)}\ 7$


 

	\item (\textit{AHSME 1957}) A man makes a trip by automobile at an average speed of 50 mph. He returns over the same route at an average speed of $45$ mph. 
His average speed for the entire trip is: 


 $\textbf{(A)}\ 47\frac{7}{19}\qquad  \textbf{(B)}\ 47\frac{1}{4}\qquad  \textbf{(C)}\ 47\frac{1}{2}\qquad  \textbf{(D)}\ 47\frac{11}{19}\qquad \textbf{(E)}\ \text{none of these}$


 

	\item (\textit{AHSME 1957}) Start with the theorem "If two angles of a triangle are equal, the triangle is isosceles," and the following four statements: 


 1. If two angles of a triangle are not equal, the triangle is not isosceles. 


 2. The base angles of an isosceles triangle are equal. 


 3. If a triangle is not isosceles, then two of its angles are not equal. 


 4. A necessary condition that two angles of a triangle be equal is that the triangle be isosceles. 


 Which combination of statements contains only those which are logically equivalent to the given theorem? 


 $\textbf{(A)}\ 1,\,2,\,3,\,4 \qquad \textbf{(B)}\ 1,\,2,\,3\qquad \textbf{(C)}\ 2,\,3,\,4\qquad \textbf{(D)}\ 1,\,2\qquad\textbf{(E)}\ 3,\,4$


 

	\item (\textit{AHSME 1957}) If $\sqrt{x - 1} - \sqrt{x + 1} + 1 = 0$, then $4x$ equals: 


 $\textbf{(A)}\ 5 \qquad \textbf{(B)}\ 4\sqrt{-1}\qquad \textbf{(C)}\ 0\qquad \textbf{(D)}\ 1\frac{1}{4}\qquad \textbf{(E)}\ \text{no real value}$


 

	\item (\textit{AHSME 1957}) The graph of $x^2 + y = 10$ and the graph of $x + y = 10$ meet in two points. The distance between these two points is: 


 $\textbf{(A)}\ \text{less than 1} \qquad \textbf{(B)}\ 1\qquad \textbf{(C)}\ \sqrt{2}\qquad \textbf{(D)}\ 2\qquad\textbf{(E)}\ \text{more than 2}$


 

	\item (\textit{AHSME 1957}) If the square of a number of two digits is decreased by the square of the number formed by reversing the digits, 
then the result is not always divisible by: 


 $\textbf{(A)}\ 9 \qquad  \textbf{(B)}\ \text{the product of the digits}\qquad  \textbf{(C)}\ \text{the sum of the digits}\qquad \textbf{(D)}\ \text{the difference of the digits}\qquad \textbf{(E)}\ 11$


 

	\item (\textit{AHSME 1957}) The vertices of $\triangle PQR$ have coordinates as follows: $P(0,a),\,Q(b,0),\,R(c,d)$, where $a,\,b,\,c$ and $d$ are positive. 
The origin and point $R$ lie on opposite sides of $PQ$. The area of $\triangle PQR$ may be found from the expression: 


 $\textbf{(A)}\ \frac{ab + ac + bc + cd}{2} \qquad  \textbf{(B)}\ \frac{ac + bd - ab}{2}\qquad  \textbf{(C)}\ \frac{ab-ac-bd}{2}\qquad \textbf{(D)}\ \frac{ac+bd+ab}{2}\qquad \textbf{(E)}\ \frac{ac+bd-ab-cd}{2}$


 

	\item (\textit{AHSME 1957}) From a point within a triangle, line segments are drawn to the vertices. 
A necessary and sufficient condition that the three triangles thus formed have equal areas is that the point be: 


 $\textbf{(A)}\ \text{the center of the inscribed circle} \qquad \\  \textbf{(B)}\ \text{the center of the circumscribed circle}\qquad\\  \textbf{(C)}\ \text{such that the three angles formed at the point each be }{120^\circ}\qquad\\  \textbf{(D)}\ \text{the intersection of the altitudes of the triangle}\qquad\\  \textbf{(E)}\ \text{the intersection of the medians of the triangle}$


 

	\item (\textit{AHSME 1957}) The sum of the reciprocals of the roots of the equation $x^2 + px + q = 0$ is: 


 $\textbf{(A)}\ -\frac{p}{q} \qquad \textbf{(B)}\ \frac{q}{p}\qquad \textbf{(C)}\ \frac{p}{q}\qquad \textbf{(D)}\ -\frac{q}{p}\qquad\textbf{(E)}\ pq$


 

	\item (\textit{AHSME 1957}) If $a$ and $b$ are positive and $a\not= 1,\,b\not= 1$, then the value of $b^{\log_b{a}}$ is: 


 $\textbf{(A)}\ \text{dependent upon }{b} \qquad \textbf{(B)}\ \text{dependent upon }{a}\qquad \textbf{(C)}\ \text{dependent upon }{a}\text{ and }{b}\qquad\textbf{(D)}\ \text{zero}\qquad\textbf{(E)}\ \text{one}$


 

	\item (\textit{AHSME 1957}) The relation $x^2(x^2 - 1)\ge 0$ is true only for: 


 $\textbf{(A)}\ x \ge 1\qquad \textbf{(B)}\ - 1 \le x \le 1\qquad \textbf{(C)}\ x = 0,\, x = 1,\, x = - 1\qquad \\\textbf{(D)}\ x = 0,\, x\le-1,\, x\ge 1\qquad\textbf{(E)}\ x\ge 0$


 

	\item (\textit{AHSME 1957}) The sum of the squares of the first n positive integers is given by the expression $\frac{n(n + c)(2n + k)}{6}$, 
if $c$ and $k$ are, respectively: 


 $\textbf{(A)}\ {1}\text{ and }{2} \qquad \textbf{(B)}\ {3}\text{ and }{5}\qquad \textbf{(C)}\ {2}\text{ and }{2}\qquad\textbf{(D)}\ {1}\text{ and }{1}\qquad\textbf{(E)}\ {2}\text{ and }{1}$


 

	\item (\textit{AHSME 1957}) A regular octagon is to be formed by cutting equal isosceles right triangles from the corners of a square. 
If the square has sides of one unit, the leg of each of the triangles has length: 


 $\textbf{(A)}\ \frac{2 + \sqrt{2}}{3} \qquad \textbf{(B)}\ \frac{2 - \sqrt{2}}{2}\qquad \textbf{(C)}\ \frac{1+\sqrt{2}}{2}\qquad\textbf{(D)}\ \frac{1+\sqrt{2}}{3}\qquad\textbf{(E)}\ \frac{2-\sqrt{2}}{3}$


 

	\item (\textit{AHSME 1957}) The largest of the following integers which divides each of the numbers of the sequence $1^5 - 1,\, 2^5 - 2,\, 3^5 - 3,\, \cdots, n^5 - n, \cdots$ is: 


 $\textbf{(A)}\ 1 \qquad \textbf{(B)}\ 60 \qquad \textbf{(C)}\ 15 \qquad \textbf{(D)}\ 120\qquad \textbf{(E)}\ 30$


 

	\item (\textit{AHSME 1957}) If $9^{x + 2} = 240 + 9^x$, then the value of $x$ is: 


 $\textbf{(A)}\ 0.1 \qquad \textbf{(B)}\ 0.2\qquad \textbf{(C)}\ 0.3\qquad \textbf{(D)}\ 0.4\qquad \textbf{(E)}\ 0.5$


 

	\item (\textit{AHSME 1957}) The points that satisfy the system $x + y = 1,\, x^2 + y^2 < 25$, constitute the following set: 


 $\textbf{(A)}\ \text{only two points} \qquad \\  \textbf{(B)}\ \text{an arc of a circle}\qquad \\  \textbf{(C)}\ \text{a straight line segment not including the end-points}\qquad\\  \textbf{(D)}\ \text{a straight line segment including the end-points}\qquad\\  \textbf{(E)}\ \text{a single point}$


 

	\item (\textit{AHSME 1957}) Side $AC$ of right triangle $ABC$ is divided into $8$ equal parts. Seven line segments parallel to $BC$ are drawn to $AB$ from the points of division. 
If $BC = 10$, then the sum of the lengths of the seven line segments: 


 $\textbf{(A)}\ \text{cannot be found from the given information} \qquad \textbf{(B)}\ \text{is }{33}\qquad \textbf{(C)}\ \text{is }{34}\qquad\textbf{(D)}\ \text{is }{35}\qquad\textbf{(E)}\ \text{is }{45}$


 

	\item (\textit{AHSME 1957}) If $x + y = 1$, then the largest value of $xy$ is: 


 $\textbf{(A)}\ 1\qquad \textbf{(B)}\ 0.5\qquad \textbf{(C)}\ \text{an irrational number about }{0.4}\qquad \textbf{(D)}\ 0.25\qquad\textbf{(E)}\ 0$


 

	\item (\textit{AHSME 1957}) In right triangle $ABC, BC = 5, AC = 12$, and $AM = x; \overline{MN} \perp \overline{AC}, \overline{NP} \perp \overline{BC}$; 
$N$ is on $AB$. If $y = MN + NP$, one-half the perimeter of rectangle $MCPN$, then: 


 
                    \begin{figure}[H]
                        \centering
                        \includegraphics[width=0.5\textwidth]{Problem/AHSME_I_1957_37_Problem_Figure_1.png}
                    \end{figure}
                    


 $\textbf{(A)}\ y = \frac {1}{2}(5 + 12) \qquad  \textbf{(B)}\ y = \frac {5x}{12} + \frac {12}{5}\qquad  \textbf{(C)}\ y =\frac{144-7x}{12}\qquad \\ \textbf{(D)}\ y = 12\qquad \qquad\quad\,\,  \textbf{(E)}\ y = \frac {5x}{12} + 6$


 

	\item (\textit{AHSME 1957}) From a two-digit number $N$ we subtract the number with the digits reversed and find that the result is a positive perfect cube. Then: 


 $\textbf{(A)}\ {N}\text{ cannot end in 5}\qquad\\  \textbf{(B)}\ {N}\text{ can end in any digit other than 5}\qquad \\  \textbf{(C)}\ {N}\text{ does not exist}\qquad\\  \textbf{(D)}\ \text{there are exactly 7 values for }{N}\qquad\\  \textbf{(E)}\ \text{there are exactly 10 values for }{N}$


 

	\item (\textit{AHSME 1957}) Two men set out at the same time to walk towards each other from $M$ and $N$, $72$ miles apart. 
The first man walks at the rate of $4$ mph. The second man walks $2$ miles the first hour, $2\frac {1}{2}$ 
miles the second hour, $3$ miles the third hour, and so on in arithmetic progression. Then the men will meet: 


 $\textbf{(A)}\ \text{in 7 hours} \qquad  \textbf{(B)}\ \text{in }{8\frac {1}{4}}\text{ hours}\qquad  \textbf{(C)}\ \text{nearer }{M}\text{ than }{N}\qquad\\  \textbf{(D)}\ \text{nearer }{N}\text{ than }{M}\qquad \textbf{(E)}\ \text{midway between }{M}\text{ and }{N}$


 

	\item (\textit{AHSME 1957}) If the parabola $y = -x^2 + bx - 8$ has its vertex on the $x$-axis, then $b$ must be: 


 $\textbf{(A)}\ \text{a positive integer}\qquad \\  \textbf{(B)}\ \text{a positive or a negative rational number}\qquad \\  \textbf{(C)}\ \text{a positive rational number}\qquad\\  \textbf{(D)}\ \text{a positive or a negative irrational number}\qquad\\  \textbf{(E)}\ \text{a negative irrational number}$


 

	\item (\textit{AHSME 1957}) Given the system of equations 

 \[ax + (a - 1)y = 1\]

 \[(a + 1)x - ay = 1\]
For which one of the following values of $a$ is there no solution $x$ and $y$? 


 $\textbf{(A)}\ 1\qquad  \textbf{(B)}\ 0\qquad  \textbf{(C)}\ - 1\qquad  \textbf{(D)}\ \frac {\pm \sqrt {2}}{2}\qquad  \textbf{(E)}\ \pm\sqrt{2}$


 

	\item (\textit{AHSME 1957}) If $S = i^n + i^{-n}$, where $i = \sqrt{-1}$ and $n$ is an integer, then the total number of possible distinct values for $S$ is: 


 $\textbf{(A)}\ 1\qquad \textbf{(B)}\ 2\qquad \textbf{(C)}\ 3\qquad \textbf{(D)}\ 4\qquad \textbf{(E)}\ \text{more than 4}$


 

	\item (\textit{AHSME 1957}) We define a lattice point as a point whose coordinates are integers, zero admitted. 
Then the number of lattice points on the boundary and inside the region bounded by the $x$-axis, 
the line $x = 4$, and the parabola $y = x^2$ is: 


 $\textbf{(A)}\ 24 \qquad  \textbf{(B)}\ 35\qquad  \textbf{(C)}\ 34\qquad  \textbf{(D)}\ 30\qquad  \textbf{(E)}\ \infty$


 

	\item (\textit{AHSME 1957}) In $\triangle ABC, AC = CD$ and $\angle CAB - \angle ABC = 30^\circ$. Then $\angle BAD$ is: 


 
                    \begin{figure}[H]
                        \centering
                        \includegraphics[width=0.5\textwidth]{Problem/AHSME_I_1957_44_Problem_Figure_1.png}
                    \end{figure}
                    


 $\textbf{(A)}\ 30^\circ\qquad\textbf{(B)}\ 20^\circ\qquad\textbf{(C)}\ 22\frac{1}{2}^\circ\qquad\textbf{(D)}\ 10^\circ\qquad\textbf{(E)}\ 15^\circ$


 

	\item (\textit{AHSME 1957}) If two real numbers $x$ and $y$ satisfy the equation $\frac{x}{y} = x - y$, then: 


 $\textbf{(A)}\ {x \ge 4}\text{ or }{x \le 0}\qquad \\  \textbf{(B)}\ {y}\text{ can equal }{1}\qquad \\  \textbf{(C)}\ \text{both }{x}\text{ and }{y}\text{ must be irrational}\qquad\\  \textbf{(D)}\ {x}\text{ and }{y}\text{ cannot both be integers}\qquad\\  \textbf{(E)}\ \text{both }{x}\text{ and }{y}\text{ must be rational}$


 

	\item (\textit{AHSME 1957}) Two perpendicular chords intersect in a circle. The segments of one chord are $3$ and $4$; 
the segments of the other are $6$ and $2$. Then the diameter of the circle is: 


 $\textbf{(A)}\ \sqrt{89}\qquad  \textbf{(B)}\ \sqrt{56}\qquad  \textbf{(C)}\ \sqrt{61}\qquad  \textbf{(D)}\ \sqrt{75}\qquad \textbf{(E)}\ \sqrt{65}$


 

	\item (\textit{AHSME 1957}) In circle $O$, the midpoint of radius $OX$ is $Q$; at $Q$, $\overline{AB} \perp \overline{XY}$. The semi-circle with $\overline{AB}$ 
as diameter intersects $\overline{XY}$ in $M$. Line $\overline{AM}$ intersects circle $O$ in $C$, and line $\overline{BM}$ 
intersects circle $O$ in $D$. Line $\overline{AD}$ is drawn. Then, if the radius of circle $O$ is $r$, $AD$ is: 


 
                    \begin{figure}[H]
                        \centering
                        \includegraphics[width=0.5\textwidth]{Problem/AHSME_I_1957_47_Problem_Figure_1.png}
                    \end{figure}
                    


 $\textbf{(A)}\ r\sqrt{2}\qquad\textbf{(B)}\ r\qquad\textbf{(C)}\ \text{not a side of an inscribed regular polygon}\qquad\textbf{(D)}\ \frac{r\sqrt{3}}{2}\qquad\textbf{(E)}\ r\sqrt{3}$


 

	\item (\textit{AHSME 1957}) Let $ABC$ be an equilateral triangle inscribed in circle $O$. $M$ is a point on arc $BC$. 
Lines $\overline{AM}$, $\overline{BM}$, and $\overline{CM}$ are drawn. Then $AM$ is: 


 
                    \begin{figure}[H]
                        \centering
                        \includegraphics[width=0.5\textwidth]{Problem/AHSME_I_1957_48_Problem_Figure_1.png}
                    \end{figure}
                    


 $\textbf{(A)}\ \text{equal to }{BM + CM}\qquad  \textbf{(B)}\ \text{less than }{BM + CM}\qquad \\ \textbf{(C)}\ \text{greater than }{BM+CM}\qquad \\ \textbf{(D)}\ \text{equal, less than, or greater than }{BM + CM}\text{, depending upon the position of } {M}\qquad  \\ \textbf{(E)}\ \text{none of these}$


 

	\item (\textit{AHSME 1957}) The parallel sides of a trapezoid are $3$ and $9$. The non-parallel sides are $4$ and $6$. 
A line parallel to the bases divides the trapezoid into two trapezoids of equal perimeters. 
The ratio in which each of the non-parallel sides is divided is: 


 
                    \begin{figure}[H]
                        \centering
                        \includegraphics[width=0.5\textwidth]{Problem/AHSME_I_1957_49_Problem_Figure_1.png}
                    \end{figure}
                    


 $\textbf{(A)}\ 4: 3\qquad\textbf{(B)}\ 3: 2\qquad\textbf{(C)}\ 4: 1\qquad\textbf{(D)}\ 3: 1\qquad\textbf{(E)}\ 6: 1$


 

	\item (\textit{AHSME 1957}) In circle $O$, $G$ is a moving point on diameter $\overline{AB}$. $\overline{AA'}$ is drawn perpendicular to $\overline{AB}$ and equal to $\overline{AG}$. 
$\overline{BB'}$ is drawn perpendicular to $\overline{AB}$, on the same side of diameter $\overline{AB}$ as $\overline{AA'}$, 
and equal to $BG$. Let $O'$ be the midpoint of $\overline{A'B'}$. Then, as $G$ moves from $A$ to $B$, point $O'$: 


 $\textbf{(A)}\ \text{moves on a straight line parallel to }{AB}\qquad \\ \textbf{(B)}\ \text{remains stationary}\qquad\\  \textbf{(C)}\ \text{moves on a straight line perpendicular to }{AB}\qquad\\  \textbf{(D)}\ \text{moves in a small circle intersecting the given circle}\qquad\\  \textbf{(E)}\ \text{follows a path which is neither a circle nor a straight line}$


 

\end{enumerate}